\section{Summary}

Value of Information (VOI) analysis asks which uncertainty matters to a
decision, what evidence could reduce it, and whether the expected improvement
justifies the cost of learning. In health technology assessment this may mean
choosing a trial, targeting a subgroup, validating a prediction model, or
improving uptake. The formulation can also be used for environmental
monitoring, engineering reliability, manufacturing quality, energy planning,
and investment decisions when uncertain consequences are mapped to a common
decision-value scale.

The difficult part is usually not the textbook expected value of perfect
information (EVPI) equation. Analysts must decide what counts as a strategy,
whose outcomes matter, which population and time horizon to scale to, whether a
model form is credible, whether evidence transports to the target setting, and
whether a recommendation will be implemented. Different choices can change
the preferred action even when the same simulation output is used. The
\texttt{voiage} package records these choices in the analysis inputs,
diagnostics, and returned result rather than leaving them as unrecorded
preprocessing.

The implemented elements described in this paper are:
\begin{itemize}
  \item a stable analysis path from probabilistic net benefits to
  EVPI, expected value of partial perfect information (EVPPI), expected value
  of sample information (EVSI), expected net benefit of sampling (ENBS),
  dominance, the cost-effectiveness acceptability frontier (CEAF),
  heterogeneity, and structural VOI\@;
  \item a labelled decision representation that records sample, strategy,
  parameter, subgroup, perspective, population, and study-design semantics;
  \item a maturity-labelled frontier that represents value questions
  about equity, perspective, implementation, validation, transportability,
  data quality, computation, and staged decisions; and
  \item diagnostic result envelopes that retain method, seed, warning,
  scaling, maturity, and provenance information that supports interpretation
  of a value estimate.
\end{itemize}

Version 1.0.0 routes selected stable aggregation kernels through
binding-independent Rust components. Python retains model orchestration,
validation for broader method paths, labelled data, plots, and
reports~\cite{voiage2026}. The R and Julia bindings currently expose direct
native EVPI only; wider application programming interface (API) parity is not
claimed. Definitions and abbreviations used throughout the paper are collected in
Sections~\ref{sec:glossary} and~\ref{sec:abbreviations}.
